% !TEX program = xelatex
% Options for packages loaded elsewhere
\PassOptionsToPackage{unicode}{hyperref}
\PassOptionsToPackage{hyphens}{url}
\PassOptionsToPackage{dvipsnames,svgnames,x11names}{xcolor}
\documentclass[
  8pt,
  twocolumn]{extarticle}
\usepackage{xcolor}
\usepackage[landscape,margin=0.38in]{geometry}
\usepackage{amsmath,amssymb}
\setcounter{secnumdepth}{-\maxdimen} % remove section numbering
\usepackage{iftex}
\ifPDFTeX
  \usepackage[T1]{fontenc}
  \usepackage[utf8]{inputenc}
  \usepackage{textcomp} % provide euro and other symbols
\else % if luatex or xetex
  \usepackage{unicode-math} % this also loads fontspec
  \defaultfontfeatures{Scale=MatchLowercase}
  \defaultfontfeatures[\rmfamily]{Ligatures=TeX,Scale=1}
\fi
\usepackage{lmodern}
\ifPDFTeX\else
  % xetex/luatex font selection
\fi
% Use upquote if available, for straight quotes in verbatim environments
\IfFileExists{upquote.sty}{\usepackage{upquote}}{}
\IfFileExists{microtype.sty}{% use microtype if available
  \usepackage[]{microtype}
  \UseMicrotypeSet[protrusion]{basicmath} % disable protrusion for tt fonts
}{}
\makeatletter
\@ifundefined{KOMAClassName}{% if non-KOMA class
  \IfFileExists{parskip.sty}{%
    \usepackage{parskip}
  }{% else
    \setlength{\parindent}{0pt}
    \setlength{\parskip}{6pt plus 2pt minus 1pt}}
}{% if KOMA class
  \KOMAoptions{parskip=half}}
\makeatother
\setlength{\emergencystretch}{3em} % prevent overfull lines
\providecommand{\tightlist}{%
  \setlength{\itemsep}{0pt}\setlength{\parskip}{0pt}}
\usepackage{microtype}
\usepackage{mathtools}
\usepackage{amssymb}
\setlength{\columnsep}{0.32in}
\setlength{\parindent}{0pt}
\setlength{\parskip}{2pt}
\setlength{\abovedisplayskip}{4pt}
\setlength{\belowdisplayskip}{4pt}
\setlength{\abovedisplayshortskip}{3pt}
\setlength{\belowdisplayshortskip}{3pt}
\allowdisplaybreaks
\sloppy
\emergencystretch=3em
\usepackage{bookmark}
\IfFileExists{xurl.sty}{\usepackage{xurl}}{} % add URL line breaks if available
\urlstyle{same}
\hypersetup{
  colorlinks=true,
  linkcolor={black},
  filecolor={Maroon},
  citecolor={Blue},
  urlcolor={black},
  pdfcreator={LaTeX via pandoc}}

\author{}
\date{}

\begin{document}

\section{Formula Sheet Part 1: Foundations, Co-Movement, and
Portfolios}\label{formula-sheet-part-1-foundations-co-movement-and-portfolios}

This is a two-column topic split from
\texttt{All\_Formulas\_From\_Chatnotes.md}.

\subsection{Conventions}\label{conventions}

\[
t,T,h,k,n,p,q,N
\]

Variables: \(t\) is a time index; \(T\) is sample size or the final
observed time; \(h\) is a forecast horizon; \(k\) is a lag or number of
model parameters depending on context; \(n\) is a return horizon;
\(p,q\) are lag orders; \(N\) is the number of assets.

\[
F_t
\]

Variables: \(F_t\) is the information set available at time \(t\).

Returns may be written as decimals or percentages. Check the scale
before applying VaR or standard deviation formulas.

\subsection{Prices, Returns, and VaR}\label{prices-returns-and-var}

\subsubsection{Simple return}\label{simple-return}

\[
R_t=\frac{P_t-P_{t-1}}{P_{t-1}}=\frac{P_t}{P_{t-1}}-1
\]

Variables: \(R_t\) is the simple return from \(t-1\) to \(t\); \(P_t\)
is the price at time \(t\); \(P_{t-1}\) is the previous price.

\subsubsection{Gross return}\label{gross-return}

\[
1+R_t=\frac{P_t}{P_{t-1}}
\]

Variables: \(1+R_t\) is the gross return, the multiplier applied to
wealth.

\subsubsection{Return with dividends}\label{return-with-dividends}

\[
R_t=\frac{P_t+D_t-P_{t-1}}{P_{t-1}}
\]

Variables: \(D_t\) is the dividend paid during period \(t\).

\subsubsection{Log return}\label{log-return}

\[
r_t=\log(1+R_t)=\log(P_t)-\log(P_{t-1})=\log\left(\frac{P_t}{P_{t-1}}\right)
\]

Variables: \(r_t\) is the log return; \(\log\) is the natural logarithm.

\subsubsection{Small-return
approximation}\label{small-return-approximation}

\[
\log(1+R_t)\approx R_t
\]

Variables: \(R_t\) is the simple return; the approximation works best
when \(R_t\) is small.

\subsubsection{Multi-period gross
return}\label{multi-period-gross-return}

\[
1+R_t(k)=\prod_{j=0}^{k-1}(1+R_{t-j})
\]

Variables: \(R_t(k)\) is the \(k\)-period simple return ending at \(t\);
\(j\) indexes the returns being multiplied.

\subsubsection{Multi-period simple
return}\label{multi-period-simple-return}

\[
R_t(k)=\prod_{j=0}^{k-1}(1+R_{t-j})-1
\]

Variables: \(R_t(k)\) is the total simple return over \(k\) periods.

\subsubsection{Equity curve for one
dollar}\label{equity-curve-for-one-dollar}

\[
V_t(k)=\prod_{j=0}^{k-1}(1+R_{t-j})
\]

Variables: \(V_t(k)\) is the value at time \(t\) of one dollar invested
over \(k\) periods.

\subsubsection{Wealth growth}\label{wealth-growth}

\[
V_T=V_0\prod_{t=1}^{T}(1+R_t)
\]

Variables: \(V_0\) is starting wealth; \(V_T\) is ending wealth after
\(T\) periods.

\subsubsection{Multi-period log return}\label{multi-period-log-return}

\[
r_t(k)=\sum_{j=0}^{k-1}r_{t-j}=r_t+r_{t-1}+\cdots+r_{t-k+1}
\]

Variables: \(r_t(k)\) is the \(k\)-period log return.

\subsubsection{Total log return between two
prices}\label{total-log-return-between-two-prices}

\[
r(1,T)=\sum_{t=1}^{T}r_t=\log\left(\frac{P_T}{P_0}\right)
\]

Variables: \(P_0\) is the starting price; \(P_T\) is the ending price.

\subsubsection{Index log return}\label{index-log-return}

\[
r_t=\log(Index_t)-\log(Index_{t-1})
\]

Variables: \(Index_t\) is the index level at time \(t\).

\subsubsection{Excess returns}\label{excess-returns}

\[
\begin{aligned}
\text{asset excess return} &= \text{asset return}-R_f \\
\text{market excess return} &= \text{market return}-R_f
\end{aligned}
\]

Variables: \(R_f\) is the risk-free return.

\subsubsection{Market beta
interpretation}\label{market-beta-interpretation}

\[
\begin{aligned}
\beta>1 &\Rightarrow \text{aggressive, more sensitive than the market} \\
\beta=1 &\Rightarrow \text{tracks the market} \\
0<\beta<1 &\Rightarrow \text{defensive, less sensitive than the market} \\
\beta=0 &\Rightarrow \text{unrelated to market movements}
\end{aligned}
\]

Variables: \(\beta\) is market sensitivity.

\subsubsection{Market capitalisation}\label{market-capitalisation}

\[
\text{market capitalisation}=\text{share price}\times\text{number of shares outstanding}
\]

Variables: market capitalisation is firm size measured using market
value.

\subsubsection{Book-to-market ratio}\label{book-to-market-ratio}

\[
\operatorname{BtM}=\frac{\text{book value of equity}}{\text{market value of equity}}
\]

Variables: \(\operatorname{BtM}\) is the book-to-market ratio.

\subsubsection{Normal return model}\label{normal-return-model}

\[
R\sim\mathcal N(\mu,\sigma^2)
\]

Variables: \(R\) is a return; \(\mu\) is mean return; \(\sigma^2\) is
variance; \(\sigma\) is standard deviation.

\subsubsection{Normal quantile}\label{normal-quantile}

\[
q_\alpha=\mu+\sigma z_\alpha
\]

Variables: \(q_\alpha\) is the lower-tail return quantile; \(\alpha\) is
the tail probability; \(z_\alpha\) is the standard normal cutoff.

\subsubsection{Common normal cutoffs}\label{common-normal-cutoffs}

\[
\begin{aligned}
z_{0.05}&=-1.645\\
z_{0.025}&=-1.96
\end{aligned}
\]

Variables: \(z_{0.05}\) is the 5 percent lower-tail cutoff;
\(z_{0.025}\) is the 2.5 percent lower-tail cutoff.

\subsubsection{One-asset VaR}\label{one-asset-var}

\[
\operatorname{VaR}_\alpha=|W_0q_\alpha|
\]

Variables: \(\operatorname{VaR}_\alpha\) is Value at Risk at tail
probability \(\alpha\); \(W_0\) is the dollar value invested.

\subsubsection{Zero-mean 5 percent normal
VaR}\label{zero-mean-5-percent-normal-var}

\[
\begin{aligned}
q_{0.05}&=-1.645\sigma\\
\operatorname{VaR}_{0.05}&=|W_0(-1.645\sigma)|
\end{aligned}
\]

Variables: \(q_{0.05}\) is the 5 percent return cutoff; \(\sigma\) is
return volatility.

\subsection{Random Variables and Distribution
Properties}\label{random-variables-and-distribution-properties}

\subsubsection{Probability under a
density}\label{probability-under-a-density}

\[
\Pr(a\le X\le b)=\text{area under }f(x)\text{ from }a\text{ to }b
\]

Variables: \(X\) is a random variable; \(a,b\) are cutoffs; \(f(x)\) is
the probability density.

\subsubsection{Normal distribution
notation}\label{normal-distribution-notation}

\[
X\sim\mathcal N(\mu,\sigma^2)
\]

Variables: \(X\) is a random variable; \(\mu\) is the mean; \(\sigma^2\)
is the variance.

\subsubsection{Expected value}\label{expected-value}

\[
\mathbb E[X]=\mu
\]

Variables: \(\mathbb E[\cdot]\) means expected value.

\subsubsection{State-contingent expected
return}\label{state-contingent-expected-return}

\[
\mathbb E(R)=\sum_s R_s\Pr(S=s)
\]

Variables: \(R_s\) is the return in state \(s\); \(S\) is the state
variable.

\subsubsection{Expected value linearity}\label{expected-value-linearity}

\[
\begin{aligned}
\mathbb E(a)&=a\\
\mathbb E(aX)&=a\mathbb E(X)\\
\mathbb E(a+X)&=a+\mathbb E(X)\\
\mathbb E(X+Y)&=\mathbb E(X)+\mathbb E(Y)
\end{aligned}
\]

Variables: \(a\) is a constant; \(X,Y\) are random variables.

\subsubsection{Variance}\label{variance}

\[
\operatorname{Var}(X)=\mathbb E[(X-\mu)^2]
\]

Variables: \(\operatorname{Var}(X)\) is the variance of \(X\).

\subsubsection{Standard deviation}\label{standard-deviation}

\[
\operatorname{sd}(X)=\sqrt{\operatorname{Var}(X)}
\]

Variables: \(\operatorname{sd}(X)\) is the standard deviation of \(X\).

\subsubsection{Variance rules}\label{variance-rules}

\[
\begin{aligned}
\operatorname{Var}(a)&=0\\
\operatorname{Var}(aX)&=a^2\operatorname{Var}(X)\\
\operatorname{Var}(a+X)&=\operatorname{Var}(X)
\end{aligned}
\]

Variables: \(a\) is a constant.

\subsubsection{Variance of a sum}\label{variance-of-a-sum}

\[
\operatorname{Var}(X+Y)=\operatorname{Var}(X)+\operatorname{Var}(Y)+2\operatorname{Cov}(X,Y)
\]

Variables: \(\operatorname{Cov}(X,Y)\) is covariance between \(X\) and
\(Y\).

\subsubsection{Standardising a normal
variable}\label{standardising-a-normal-variable}

\[
Z=\frac{X-\mu}{\sigma}\sim\mathcal N(0,1)
\]

Variables: \(Z\) is a standard normal variable; \(\sigma\) is the
standard deviation of \(X\).

\subsubsection{Probability conversion after
standardising}\label{probability-conversion-after-standardising}

\[
\Pr(X>a)=\Pr\left(\frac{X-\mu}{\sigma}>\frac{a-\mu}{\sigma}\right)=\Pr(Z>z)
\]

Variables: \(z=(a-\mu)/\sigma\) is the standardised cutoff.

\subsubsection{Lower-tail normal
probability}\label{lower-tail-normal-probability}

\[
\Pr(r_t<x)=\Pr\left(Z<\frac{x-\mu}{\sigma}\right)
\]

Variables: \(r_t\) is a return; \(x\) is a cutoff.

\subsubsection{Return scale convention}\label{return-scale-convention}

\[
\begin{aligned}
0.01&=1\text{ percent, if returns are decimals}\\
1&=1\text{ percent, if returns are stored as percentage points}
\end{aligned}
\]

Variables: the same economic return can be represented on different
numeric scales.

\subsubsection{Sample mean}\label{sample-mean}

\[
\bar r=\frac{1}{T}\sum_{t=1}^{T}r_t
\]

Variables: \(\bar r\) is the sample average return.

\subsubsection{Sample variance and standard
deviation}\label{sample-variance-and-standard-deviation}

\[
s^2=\frac{1}{T-1}\sum_{t=1}^{T}(r_t-\bar r)^2,\qquad s=\sqrt{s^2}
\]

Variables: \(s^2\) is sample variance; \(s\) is sample standard
deviation.

\subsubsection{Skewness and kurtosis}\label{skewness-and-kurtosis}

\[
\text{skewness}=\frac{\mathbb E[(r_t-\mu)^3]}{\sigma^3},\qquad
\text{kurtosis}=\frac{\mathbb E[(r_t-\mu)^4]}{\sigma^4}
\]

Variables: skewness measures asymmetry; kurtosis measures tail
thickness.

\subsection{Co-Movement, Conditional Moments, and
Factors}\label{co-movement-conditional-moments-and-factors}

\subsubsection{Joint probability}\label{joint-probability}

\[
\Pr(X=x,Y=y)
\]

Variables: \(X,Y\) are random variables; \(x,y\) are outcomes.

\subsubsection{Marginal probability}\label{marginal-probability}

\[
\Pr(X=x)=\sum_y\Pr(X=x,Y=y)
\]

Variables: the sum is over all possible values of \(Y\).

\subsubsection{Conditional probability}\label{conditional-probability}

\[
\Pr(X=x\mid Y=y)=\frac{\Pr(X=x,Y=y)}{\Pr(Y=y)}
\]

Variables: \(\Pr(X=x\mid Y=y)\) is the probability of \(X=x\) given
\(Y=y\).

\subsubsection{Independence}\label{independence}

\[
\Pr(X=x\mid Y=y)=\Pr(X=x)
\]

Variables: independence means knowing \(Y\) does not change the
probability distribution of \(X\).

\subsubsection{Conditional expectation}\label{conditional-expectation}

\[
\mathbb E[X\mid Y]
\]

Variables: \(\mathbb E[X\mid Y]\) is the expected value of \(X\)
conditional on \(Y\).

\subsubsection{Law of iterated
expectations}\label{law-of-iterated-expectations}

\[
\mathbb E[X]=\mathbb E[\mathbb E[X\mid Y]]
\]

Variables: \(Y\) is the conditioning variable.

\subsubsection{Conditional variance}\label{conditional-variance}

\[
\operatorname{Var}(Y\mid X)=\mathbb E[(Y-\mathbb E[Y\mid X])^2\mid X]
\]

Variables: \(\operatorname{Var}(Y\mid X)\) is remaining uncertainty
about \(Y\) after conditioning on \(X\).

\subsubsection{Law of total variance}\label{law-of-total-variance}

\[
\operatorname{Var}(Y)=\mathbb E[\operatorname{Var}(Y\mid X)]+\operatorname{Var}(\mathbb E[Y\mid X])
\]

Variables: the first term is average conditional variance; the second is
variance of conditional means.

\subsubsection{Information reduces conditional
variance}\label{information-reduces-conditional-variance}

\[
\operatorname{Var}(Y\mid \text{information})<\operatorname{Var}(Y)
\]

Variables: information is any conditioning set that improves prediction.

\subsubsection{Covariance}\label{covariance}

\[
\operatorname{Cov}(X,Y)=\mathbb E[(X-\mathbb E[X])(Y-\mathbb E[Y])]
\]

Variables: covariance measures linear co-movement.

\subsubsection{State covariance}\label{state-covariance}

\[
\operatorname{Cov}(R_A,R_B)=\sum_s p_s(R_{A,s}-\mu_A)(R_{B,s}-\mu_B)
\]

Variables: \(R_A,R_B\) are asset returns; \(p_s\) is the probability of
state \(s\); \(\mu_A,\mu_B\) are expected returns.

\subsubsection{Covariance from
correlation}\label{covariance-from-correlation}

\[
\operatorname{Cov}(R_A,R_B)=\rho_{AB}\sigma_A\sigma_B
\]

Variables: \(\rho_{AB}\) is the correlation between assets A and B;
\(\sigma_A,\sigma_B\) are standard deviations.

\subsubsection{Correlation}\label{correlation}

\[
\operatorname{Corr}(X,Y)=\frac{\operatorname{Cov}(X,Y)}{\operatorname{sd}(X)\operatorname{sd}(Y)}
\]

Variables: \(\operatorname{Corr}(X,Y)\) is the standardised covariance.

\subsubsection{Correlation from
covariance}\label{correlation-from-covariance}

\[
\rho_{AB}=\frac{\operatorname{Cov}(R_A,R_B)}{\sigma_A\sigma_B}
\]

Variables: \(\rho_{AB}\) is the correlation coefficient.

\subsubsection{Covariance sign
interpretation}\label{covariance-sign-interpretation}

\[
\begin{aligned}
\operatorname{Cov}(X,Y)>0&\Rightarrow X\text{ and }Y\text{ tend to move together}\\
\operatorname{Cov}(X,Y)<0&\Rightarrow X\text{ and }Y\text{ tend to move oppositely}\\
\operatorname{Cov}(X,Y)=0&\Rightarrow \text{no linear co-movement}
\end{aligned}
\]

Variables: the sign of covariance gives direction of linear co-movement.

\subsubsection{Group mean return}\label{group-mean-return}

\[
\bar r_g=\frac{1}{T}\sum_{t=1}^{T}r_{g,t}
\]

Variables: \(\bar r_g\) is the mean return for group \(g\); \(r_{g,t}\)
is group return at time \(t\).

\subsubsection{Generic factor long-short
return}\label{generic-factor-long-short-return}

\[
r_{\text{factor},t}=r_{\text{long},t}-r_{\text{short},t}
\]

Variables: \(r_{\text{long},t}\) is the return on the long leg;
\(r_{\text{short},t}\) is the return on the short leg.

\subsubsection{SMB factor}\label{smb-factor}

\[
SMB_t=R_{\text{small},t}-R_{\text{big},t}
\]

Variables: \(SMB_t\) is Small Minus Big at time \(t\).

\subsubsection{HML factor}\label{hml-factor}

\[
HML_t=R_{\text{high BtM},t}-R_{\text{low BtM},t}
\]

Variables: \(HML_t\) is High Minus Low book-to-market at time \(t\).

\subsubsection{Fama-French two-by-three SMB
construction}\label{fama-french-two-by-three-smb-construction}

\[
SMB=\operatorname{average}(S/L,S/N,S/H)-\operatorname{average}(B/L,B/N,B/H)
\]

Variables: \(S\) means small; \(B\) means big; \(L,N,H\) mean low,
neutral, and high book-to-market portfolios.

\subsubsection{Fama-French two-by-three HML
construction}\label{fama-french-two-by-three-hml-construction}

\[
HML=\operatorname{average}(S/H,B/H)-\operatorname{average}(S/L,B/L)
\]

Variables: \(S/H\) is the small high-BtM portfolio; \(B/H\) is the big
high-BtM portfolio; \(S/L,B/L\) are low-BtM portfolios.

\subsection{Portfolio Risk and Sharpe
Ratios}\label{portfolio-risk-and-sharpe-ratios}

\subsubsection{Two-asset portfolio
return}\label{two-asset-portfolio-return}

\[
R_p=wR_1+(1-w)R_2
\]

Variables: \(R_p\) is portfolio return; \(w\) is the weight in asset 1;
\(R_1,R_2\) are asset returns.

\subsubsection{Two-asset expected
return}\label{two-asset-expected-return}

\[
\mathbb E[R_p]=w\mathbb E[R_1]+(1-w)\mathbb E[R_2]
\]

Variables: \(\mathbb E[R_p]\) is expected portfolio return.

\subsubsection{Two-asset expected return using
means}\label{two-asset-expected-return-using-means}

\[
\mu_p=w\mu_1+(1-w)\mu_2
\]

Variables: \(\mu_p\) is portfolio mean return; \(\mu_1,\mu_2\) are asset
mean returns.

\subsubsection{Two-asset portfolio
variance}\label{two-asset-portfolio-variance}

\[
\operatorname{Var}(R_p)=w^2\sigma_1^2+(1-w)^2\sigma_2^2+2w(1-w)\sigma_{12}
\]

Variables: \(\sigma_1^2,\sigma_2^2\) are asset variances;
\(\sigma_{12}\) is covariance.

\subsubsection{Two-asset portfolio variance using explicit
weights}\label{two-asset-portfolio-variance-using-explicit-weights}

\[
\sigma_p^2=w_1^2\sigma_1^2+w_2^2\sigma_2^2+2w_1w_2\sigma_{12}
\]

Variables: \(w_1,w_2\) are portfolio weights; \(\sigma_p^2\) is
portfolio variance.

\subsubsection{Covariance from
correlation}\label{covariance-from-correlation-1}

\[
\sigma_{12}=\rho_{12}\sigma_1\sigma_2
\]

Variables: \(\rho_{12}\) is correlation between assets 1 and 2.

\subsubsection{Sharpe ratio}\label{sharpe-ratio}

\[
SR=\frac{\operatorname{mean\ return}-\text{risk-free rate}}{\text{standard deviation}}
\]

Variables: \(SR\) is the Sharpe ratio.

\subsubsection{Asset Sharpe ratio}\label{asset-sharpe-ratio}

\[
SR_i=\frac{\mu_i-r_f}{\sigma_i}
\]

Variables: \(\mu_i\) is expected return on asset \(i\); \(r_f\) is the
risk-free rate; \(\sigma_i\) is standard deviation.

\subsubsection{Portfolio Sharpe ratio}\label{portfolio-sharpe-ratio}

\[
SR_p=\frac{\mu_p-r_f}{\sigma_p}
\]

Variables: \(SR_p\) is portfolio Sharpe ratio; \(\sigma_p\) is portfolio
standard deviation.

\subsubsection{Two-asset GMV weight}\label{two-asset-gmv-weight}

\[
w_1^{GMV}=\frac{\sigma_2^2-\sigma_{12}}{\sigma_1^2+\sigma_2^2-2\sigma_{12}}
\]

Variables: \(w_1^{GMV}\) is the global minimum variance weight in asset
1.

\subsubsection{Second GMV weight}\label{second-gmv-weight}

\[
w_2^{GMV}=1-w_1^{GMV}
\]

Variables: \(w_2^{GMV}\) is the GMV weight in asset 2.

\subsubsection{Regression trick for GMV
weights}\label{regression-trick-for-gmv-weights}

\[
r_{2,t}=\beta_0+\beta_1(r_{2,t}-r_{1,t})+e_t
\]

Variables: \(r_{1,t},r_{2,t}\) are asset returns; \(\beta_0,\beta_1\)
are regression coefficients; \(e_t\) is the error.

\[
\beta_1=w_1,\qquad w_2=1-\beta_1
\]

Variables: \(\beta_1\) gives the weight in asset 1 under the regression
trick.

\subsubsection{Regression slope form used in
practice}\label{regression-slope-form-used-in-practice}

\[
b=\frac{\operatorname{Cov}(R_1,R_2)}{\operatorname{Var}(R_2)}
\]

Variables: \(b\) is a regression slope in the practice note.

\subsubsection{Many-asset GMV weights}\label{many-asset-gmv-weights}

\[
\mathbf w^{GMV}=\frac{\Sigma^{-1}\mathbf 1}{\mathbf 1'\Sigma^{-1}\mathbf 1}
\]

Variables: \(\mathbf w^{GMV}\) is the GMV weight vector; \(\Sigma\) is
the covariance matrix; \(\mathbf 1\) is a vector of ones.

\subsubsection{Portfolio return from GMV
weights}\label{portfolio-return-from-gmv-weights}

\[
r_{p,t}=w_1r_{1,t}+(1-w_1)r_{2,t}
\]

Variables: \(r_{p,t}\) is portfolio return at time \(t\).

\end{document}
